
\documentclass{article}
\usepackage{colortbl}
\usepackage{makecell}
\usepackage{multirow}
\usepackage{supertabular}

\begin{document}

\newcounter{utterance}

\centering \large Interaction Transcript for game `codenames', experiment `association\_difficult', episode 8 with reflexion{-}Qwen3.5{-}27B.
\vspace{24pt}

{ \footnotesize  \setcounter{utterance}{1}
\setlength{\tabcolsep}{0pt}
\begin{supertabular}{c@{$\;$}|p{.15\linewidth}@{}p{.15\linewidth}p{.15\linewidth}p{.15\linewidth}p{.15\linewidth}p{.15\linewidth}}
    \# & $\;$A & \multicolumn{4}{c}{Game Master} & $\;\:$B\\
    \hline

    \theutterance \stepcounter{utterance}  
    & & & \multicolumn{2}{p{0.3\linewidth}}{
        \cellcolor[rgb]{0.9,0.9,0.9}{
            \makecell[{{p{\linewidth}}}]{
                \texttt{\tiny{[GM$|$GM]}}
                {'hidden': {'team': ['head'], 'innocent': ['ash', 'sage', 'teacher', 'running', 'stomach'], 'opponent': [], 'assassin': []}, 'revealed': {'team': {'team': [], 'innocent': [], 'opponent': [], 'assassin': []}, 'opponent': {'team': [], 'innocent': [], 'opponent': [], 'assassin': []}}}
            }
        }
    }
    & & \\ \\

    \theutterance \stepcounter{utterance}  
    & & \multicolumn{4}{p{0.6\linewidth}}{
        \cellcolor[rgb]{0.9,0.9,0.9}{
            \makecell[{{p{\linewidth}}}]{
                \texttt{\tiny{[P1$\langle$GM]}}
                \texttt{You have attempted this same task before and did not succeed. These are the plans you wrote after those attempts.} \\
\texttt{Trial \#1: New plan: The previous attempt was aborted because the clue "red" was identical to one of the target words ("red"), violating the rule that the clue cannot be one of the words in the lists. Additionally, the clue "red" is too broad and risks triggering the opponent's word "boxing" (red corner) or distractors. Next time, I will select a clue that is semantically related to multiple target words but is distinct from all words in the target, opponent, distractor, and kill lists. For example, I will use "color" to link "raspberry" and "red" (and potentially "cinnamon" as a shade), ensuring the clue word itself is not present in any of the provided lists. I will strictly verify that the chosen clue is not a substring or exact match of any forbidden word before submitting.} \\
\\ 
\texttt{Trial \#2: New plan: My previous strategy failed because I prioritized grouping words by loose semantic associations (like "wood" for cinnamon and carpenter) without sufficiently evaluating the risk of triggering opponent or distractor words. Specifically, "wood" likely led to "cave" (opponent) or "stream" (distractor) due to shared environmental contexts, ending my turn prematurely. Additionally, I left high{-}value, easily guessable words like "church" and "chapel" for later, missing an opportunity to secure quick points with a safe, high{-}precision clue. Next time, I will prioritize clues that link multiple remaining target words with a very specific, low{-}ambiguity association that has no overlap with the opponent's list (e.g., using "building" or "place" for "church" and "chapel" while avoiding "cathedral" by being specific to the pair, or "fruit" for "raspberry" if safe). I will strictly cross{-}reference every potential clue against the opponent's and distractor lists to ensure no semantic bridge exists before submitting, favoring fewer, safer targets over risky, broad groupings.} \\
\\ 
\texttt{Trial \#3: New plan: My previous strategy failed because I prioritized safety and simplicity by giving one{-}word clues for single targets, which allowed the opponent to catch up and eventually win. Specifically, I missed the opportunity to group "head" and "raspberry" (or "head" and "cinnamon") under a broader, safe clue like "top" or "part" earlier in the game to clear multiple words in a single turn. Additionally, I failed to account for the high risk of the clue "body" triggering the opponent's word "hand" (as both are body parts) or the distractor "stomach," which likely caused the loss. Next time, I will aggressively group remaining words into multi{-}target clues (e.g., "top" for head/raspberry/cinnamon) to maximize points per turn, while strictly avoiding clues that share semantic categories with the opponent's remaining words (like "body" when "hand" is active). I will prioritize speed over extreme caution once the board state is clear, ensuring I clear my list before the opponent can.} \\
\\ 
\\ 
\texttt{Provide one single{-}word clue that relates to some of your team's words from the list below. You can choose to how many and to which words your clue relates to.} \\
\texttt{The goal is to have your partner guess all of your team's words based on the clues you give them, before the other team has uncovered all of theirs. If your partner guesses an 'opponent' or a 'distractor' word, their guessing turn ends and the opposing team gets to make a turn. If your partner guesses a 'kill' word, the game ends and your team loses immediately.} \\
\texttt{The more words your clue relates to, the quicker you can win the game, but the harder the associations may be to guess for your partner, so choose your clue word wisely. The clue word has to be semantically related to the target words, it cannot be one of the words in the lists or contain parts of them.} \\
\\ 
\texttt{Always give your single{-}word clue and your comma{-}separated list of related target words in the following format and make your answers as short as possible, never include any other text than is required in this form:} \\
\\ 
\texttt{CLUE: <WORD>} \\
\texttt{TARGETS: <TARGETS>} \\
\\ 
\texttt{Your team words are: head.} \\
\texttt{Your opponent's team words are: .} \\
\texttt{Distractor words are: ash, sage, teacher, running, stomach.} \\
\texttt{The kill words are: .} \\
            }
        }
    }
    & \\ \\

    \theutterance \stepcounter{utterance}  
    & \multicolumn{4}{p{0.6\linewidth}}{
        \cellcolor[rgb]{0.9,0.9,0.9}{
            \makecell[{{p{\linewidth}}}]{
                \texttt{\tiny{[P1$\rangle$GM]}}
                \texttt{CLUE: top} \\
\texttt{TARGETS: head} \\
            }
        }
    }
    & & \\ \\

    \theutterance \stepcounter{utterance}  
    & & & \multicolumn{2}{p{0.3\linewidth}}{
        \cellcolor[rgb]{0.9,0.9,0.9}{
            \makecell[{{p{\linewidth}}}]{
                \texttt{\tiny{[GM$|$GM]}}
                \texttt{top} \\
            }
        }
    }
    & & \\ \\

    \theutterance \stepcounter{utterance}  
    & & & \multicolumn{2}{p{0.3\linewidth}}{
        \cellcolor[rgb]{0.9,0.9,0.9}{
            \makecell[{{p{\linewidth}}}]{
                \texttt{\tiny{[GM$|$GM]}}
                ['head']
            }
        }
    }
    & & \\ \\

    \theutterance \stepcounter{utterance}  
    & & & \multicolumn{2}{p{0.3\linewidth}}{
        \cellcolor[rgb]{0.9,0.9,0.9}{
            \makecell[{{p{\linewidth}}}]{
                \texttt{\tiny{[GM$|$GM]}}
                {'word': 'head', 'assignment': 'team'}
            }
        }
    }
    & & \\ \\

    \theutterance \stepcounter{utterance}  
    & & & \multicolumn{2}{p{0.3\linewidth}}{
        \cellcolor[rgb]{0.9,0.9,0.9}{
            \makecell[{{p{\linewidth}}}]{
                \texttt{\tiny{[GM$|$GM]}}
                \texttt{opponent has won} \\
            }
        }
    }
    & & \\ \\

    \theutterance \stepcounter{utterance}  
    & \multicolumn{4}{p{0.6\linewidth}}{
        \cellcolor[rgb]{0.95,0.95,0.95}{
            \makecell[{{p{\linewidth}}}]{
                \texttt{\tiny{[P2$\rightarrow$P1]}}
                \texttt{Post{-}episode Reflexion (not sent to GM):} \\
\texttt{New plan: My previous attempt failed because I provided a clue ("top") that, while semantically related to "head," was too broad and likely triggered a distractor word (e.g., "running" as in "top speed" or "top of the list," or "sage" as in "top sage") or simply failed to secure the win before the opponent's turn. With only one target word ("head") and no opponent words, the optimal strategy is to use a highly specific, low{-}ambiguity clue that uniquely identifies "head" without connecting to any distractors (ash, sage, teacher, running, stomach). I will select a clue like "brain" or "cap" which has a direct, exclusive association with "head" and no plausible link to the distractors. I will strictly verify that the chosen clue does not share a semantic category or common phrase with any distractor before submitting.} \\
            }
        }
    }
    & & \\ \\

\end{supertabular}
}

\end{document}
